\section{Fixed points and centres in trees}\label{sec:trees}

\begin{setup}[Trees]\label{setup:trees}
Let $T$ be a locally finite tree and $|T|$ its geometric realisation, a
metric space in which every edge is isometric to $[0,1]$. Groups act on
$|T|$ through automorphisms of $T$; such an automorphism \emph{inverts} an
edge if it maps the edge to itself and swaps its endpoints. For a group
$G$ acting on $|T|$ let $\FixedSet{G} \subseteq |T|$ be its fixed set.
\end{setup}

\begin{lemma}[Fixed sets are convex]\label{lem:fixed-convex}
In the situation of \cref{setup:trees}, $\FixedSet{G}$ is convex: it
contains the geodesic between any two of its points.
\end{lemma}

\begin{proof}
Geodesics in $|T|$ are unique and isometries map geodesics to geodesics.
\end{proof}

\begin{lemma}[Shape of fixed sets]\label{lem:fixed-shape}
\slogan{A group acting on a tree fixes a subtree, or the midpoint of one
inverted edge, or nothing.}
In the situation of \cref{setup:trees}, exactly one of the following
holds:
\begin{enumerate}
\item $\FixedSet{G}$ contains a vertex; then $\FixedSet{G}$ is the
  realisation of the full subgraph spanned by the fixed vertices, which is
  a subtree;
\item $\FixedSet{G}$ contains no vertex and is nonempty; then
  $\FixedSet{G}$ is the midpoint of a single edge which some element of
  $G$ inverts;
\item $\FixedSet{G} = \emptyset$.
\end{enumerate}
\end{lemma}

\begin{proof}
An automorphism fixing a point in the interior of an edge maps that edge
to itself; if the point is not the midpoint, or if an endpoint is also
fixed, then both endpoints are fixed and so is the whole edge. If a
midpoint $m$ and a vertex $v$ are both fixed, the geodesic from $v$ to $m$
passes through an endpoint of the edge of $m$, which is fixed by
convexity (\cref{lem:fixed-convex}), so the edge is fixed pointwise. This
gives (1) and (2); in case (2) two distinct fixed midpoints would force a
fixed vertex between them.
\end{proof}

\begin{definition}[Centre]\label{def:centre}
Let $Y \subseteq |T|$ be nonempty, bounded and convex. Its
\emph{diameter} $\Diameter{Y}$ is the supremum of the distances between
its points, and its \emph{centre} $\Centre{Y}$ is the unique point
$c \in |T|$ with $\sup_{y \in Y} d(c, y) \leq \Diameter{Y}/2$
(\cref{lem:centre}).
\end{definition}

\begin{lemma}[Centres exist and are equivariant]\label{lem:centre}
\slogan{A bounded convex subset of a tree has a unique centre, and it is a
vertex or an edge midpoint when the set is a subtree.}
Let $Y \subseteq |T|$ be nonempty, bounded, convex and closed. Then there
is a unique point $c$ with $\sup_{y\in Y} d(c,y) \leq \Diameter{Y}/2$; it is
the midpoint of every geodesic in $Y$ of length $\Diameter{Y}$. If $Y$ is a
finite subtree, or the midpoint of an edge, then $\Centre{Y}$ is a vertex
or the midpoint of an edge. For every automorphism $\phi$ of $T$ one has
$\Centre{\phi Y} = \phi\Centre{Y}$.
\end{lemma}

\begin{proof}
Let $[a,b] \subseteq Y$ be a geodesic of length $D = \Diameter{Y}$ (it
exists since $Y$ is compact) with midpoint $m$. For $y \in Y$ let $q$ be
the point where the geodesic from $y$ to $m$ first meets $[a,b]$; if
$d(y,m) > D/2$ then $d(y,a)$ or $d(y,b)$ exceeds $D$. Conversely if
$\sup_{y} d(c, y) \leq D/2$ then $d(a,c) + d(c,b) \leq D = d(a,b)$, so $c$
lies on $[a,b]$ at distance $D/2$ from $a$, i.e.\ $c = m$. For a finite
subtree $D$ is an integer, so $m$ is a vertex or a midpoint. Equivariance
holds because the defining property is preserved by isometries.
\end{proof}

\begin{lemma}[Groups with a finite orbit fix a point]\label{lem:finite-orbit}
\slogan{A group with a finite vertex orbit fixes the centre of its convex
hull.}
In the situation of \cref{setup:trees}, if $G$ has a finite orbit $O$ on
the vertices of $T$, then $G$ fixes $\Centre{\ConvexHull{O}}$, which is a
vertex or the midpoint of an edge.
\end{lemma}

\begin{proof}
$\ConvexHull{O}$ is a finite subtree preserved by $G$; apply the
equivariance of centres (\cref{lem:centre}).
\end{proof}

\begin{lemma}[Fixed points of the standard local groups]
\label{lem:local-fixed-points}
\slogan{$\SL{2}{\pAdicIntegers{p}}$ fixes only the standard vertex, the
Iwahori subgroup only the standard edge, and its normaliser only the
midpoint.}
Let $p$ be a prime and let $\IwahoriOrder{p} := \{\smallMat{a}{b}{c}{d}
\in \MatRing{2}{\pAdicIntegers{p}} : p \mid c\}$. Acting on
$\TreeRealization{p}$:
\begin{enumerate}
\item every subgroup of $\PGL{2}{\pAdicNumbers{p}}$ containing the image of
  $\SL{2}{\pAdicIntegers{p}}$ and fixing $\LocalBaseVertex{p}$ has fixed set
  $\{\LocalBaseVertex{p}\}$;
\item every subgroup of the image of $\unit{\IwahoriOrder{p}}$ containing the
  image of $\unit{\IwahoriOrder{p}} \cap \SL{2}{\pAdicIntegers{p}}$ has fixed
  set the closed edge between $\LocalBaseVertex{p}$ and
  $\LocalStdVertex{p}{p}$;
\item every subgroup as in (2), extended by an element inverting that
  edge, has fixed set $\{\Midpoint{\LocalBaseVertex{p}}{\LocalStdVertex{p}{p}}\}$.
\end{enumerate}
\end{lemma}

\begin{proof}
The neighbours of $\LocalBaseVertex{p}$ correspond to the lines in
$\StdModule{\FF{p}}{2}$, via $L \mapsto L/p\StdModule{\pAdicIntegers{p}}{2}$;
the neighbour $\LocalStdVertex{p}{p}$ corresponds to the line spanned by
$\StdBasis{1}$. (1): $\SL{2}{\FF{p}}$ acts transitively on lines, so by
\cref{lem:fixed-shape} the fixed set contains no other vertex and no
midpoint. (2): the group fixes $\LocalBaseVertex{p}$ and $\LocalStdVertex{p}{p}$
by $\GammaZero{}$ as a fixer (\cref{lem:gamma0-fixer}). The matrices
$\smallMat{1}{b}{0}{1}$ act transitively on the $p$ lines other than
$\FF{p}\StdBasis{1}$, and the matrices $\smallMat{1}{0}{pc}{1}$ act
transitively on the $p$ neighbours of $\LocalStdVertex{p}{p}$ other than
$\LocalBaseVertex{p}$ (conjugate by a local flip). Hence no further vertex
is fixed. (3) follows from (2) and \cref{lem:fixed-shape}.
\end{proof}

\begin{lemma}[Fixed sets of principal congruence subgroups]
\label{lem:fixed-principal}
Let $p$ be a prime and $k \geq 0$. The image in $\PGL{2}{\pAdicNumbers{p}}$
of $\{g \in \SL{2}{\pAdicIntegers{p}} : g \equiv 1 \bmod p^{k}\}$ has as fixed
set the ball of radius $k$ around $\LocalBaseVertex{p}$ in
$\TreeRealization{p}$.
\end{lemma}

\begin{proof}
The vertices at distance at most $k$ from $\LocalBaseVertex{p}$ are the
classes of lattices $L$ with $p^{k}\StdModule{\pAdicIntegers{p}}{2} \subseteq L
\subseteq \StdModule{\pAdicIntegers{p}}{2}$ (up to scaling), all of which are
fixed. A vertex at distance $k+1$ corresponds to a cyclic submodule of
$\StdModule{(\ResidueRing{p^{k+1}})}{2}$ of order $p^{k+1}$, and the group
acts on these through the matrices $\equiv 1 \bmod p^{k}$, which act
nontrivially, e.g.\ $\smallMat{1}{p^{k}}{0}{1}$ or $\smallMat{1}{0}{p^{k}}{1}$
moves every such submodule other than possibly $\langle\StdBasis{1}\rangle$
respectively $\langle\StdBasis{2}\rangle$. Conclude with
\cref{lem:fixed-shape}.
\end{proof}
